\section{Extensions: theory}
\begin{frame}[plain]{Part II: Extensions}
\Large Departures from the canonical sharp RD setup
\vspace{1em}\normalsize
\begin{enumerate}\setlength{\itemsep}{.7em}
\item Local randomization
\item Fuzzy RD
\item Discrete running variables
\item Multiple cutoffs and multiple scores
\end{enumerate}
\EX{Sections 1--5}
\end{frame}

\subsection{Local randomization}
\begin{frame}{Continuity and local randomization}
\fig[.60]{experiment_rd.pdf}
\begin{items}
\item Continuity allows potential outcome means to vary with the score.
\item Local randomization explicitly models assignment in a fixed window.
\end{items}
\EX{Section 2, Figure 2}
\end{frame}

\begin{frame}{The local randomization assumptions}
Choose a window $\mathcal W=[c-w,c+w]$ around the cutoff.
\begin{items}
\item \textbf{LR1:} within $\mathcal W$, the score is unconfounded with potential outcomes and the treatment assignment mechanism is known.
\item \textbf{LR2:} within $\mathcal W$, changing the score has no direct effect on potential outcomes.
\item Both assumptions need substantive justification in the chosen window.
\item Randomness in the score alone does not rule out a direct score effect.
\end{items}
\EX{Section 2, conditions LR1--LR2}
\end{frame}

\begin{frame}{A different estimand inside the window}
\begin{columns}
\begin{column}{.54\textwidth}\fig[.64]{local_randomization.pdf}\end{column}
\begin{column}{.43\textwidth}
For fixed potential outcomes,
\[
\theta_{\mathcal W}=\frac{1}{N_{\mathcal W}}
\sum_{i:X_i\in\mathcal W}[Y_i(1)-Y_i(0)].
\]
This is a window average.\\[.7em]
The continuity target $\tau_{\mathrm{SRD}}$ is an average at $X=c$.
\end{column}
\end{columns}
\EX{Sections 2 and 2.2.2, Figure 3}
\end{frame}

\begin{frame}{The assignment mechanism}
\begin{items}
\item \textbf{Complete randomization:} condition on the observed number treated in the window. All assignments with that count have equal probability.
\item \textbf{Bernoulli assignment:} each unit receives treatment with a specified probability, allowing the treated count to vary.
\item These mechanisms imply different randomization distributions.
\item Choosing one is part of the research design and needs justification.
\end{items}
\EX{Section 2.2.1}
\end{frame}

\begin{frame}{Fisherian inference}
The sharp null specifies every unit's treatment effect:
\[
H_0^F:\ Y_i(1)=Y_i(0)\quad\text{for every }i\text{ in }\mathcal W.
\]
\begin{enumerate}\setlength{\itemsep}{.5em}
\item Under this null, impute all missing potential outcomes.
\item Reassign treatment according to the stated mechanism.
\item Recalculate a statistic, such as the difference in means.
\item Compare the observed statistic with this randomization distribution.
\end{enumerate}
\EX{Section 2.2.1}
\end{frame}

\begin{frame}{Exact tests and Monte Carlo approximations}
\begin{items}
\item Enumerating every permissible assignment gives the finite assignment distribution under the sharp null.
\item Simulation approximates that distribution when enumeration is costly.
\item A fixed seed makes the simulation reproducible. More draws reduce Monte Carlo error.
\item A displayed simulation $p$-value of zero means no simulated draw was as extreme. It is not evidence of a literally zero probability.
\end{items}
\EX{Section 2.2.1}
\end{frame}

\begin{frame}{Fisher intervals and average effect inference}
\begin{items}
\item Invert sharp hypotheses $Y_i(1)-Y_i(0)=\theta$ to form a Fisherian interval under a constant effect model.
\item A Neyman null instead concerns the average effect: $H_0^N:\theta_{\mathcal W}=0$.
\item Under complete randomization, the difference in means estimates the window average effect.
\item Large sample inference for an average effect and Fisherian inference test different hypotheses.
\end{items}
\EX{Sections 2.2.1--2.2.2}
\end{frame}

\begin{frame}{Selecting the window}
\begin{enumerate}\setlength{\itemsep}{.55em}
\item Specify a sequence of nested windows and predetermined covariates.
\item Evaluate covariate balance, beginning with the smallest window.
\item Enlarge the window while it meets the stated balance criterion.
\item Fix the rule before inspecting the main outcome results.
\end{enumerate}
\vspace{.6em}
Failing to reject imbalance does not establish randomization. If the smallest window fails, the procedure may select no window.
\EX{Section 2.2.4}
\end{frame}

\begin{frame}{Validation under local randomization}
\begin{items}
\item Examine predetermined covariates using the same window and assignment mechanism.
\item Compare counts above and below the cutoff under a justified assignment probability.
\item Use placebo cutoffs and sensitivity to nearby window choices.
\item Selecting a window and then reusing balance evidence requires caution about what the tests establish.
\end{items}
\EX{Section 2.3}
\end{frame}

\begin{frame}{When local randomization is useful}
\begin{items}
\item It can support finite sample tests when a credible small window contains few units.
\item With a continuous score, its fixed-window assumptions are typically stronger than continuity alone.
\item With a coarse discrete score, it offers an alternative to extrapolation between unobserved score values.
\item Different windows change both the units in the estimand and the plausibility of the assumptions.
\end{items}
\EX{Section 2.4}
\end{frame}

\subsection{The fuzzy RD design}
\begin{frame}{Assignment and receipt with imperfect compliance}
\fig[.57]{sharp_fuzzy.pdf}
\[
T_i=\ind\{X_i\geq c\},\qquad D_i\text{ is actual treatment received.}
\]
The discontinuity in $\Pr(D_i=1\mid X_i=x)$ is nonzero and smaller than one.
\EX{Section 3, Figure 7}
\end{frame}

\begin{frame}{Potential treatments and potential outcomes}
\[
D_i=T_iD_i(1)+(1-T_i)D_i(0).
\]
\begin{items}
\item $D_i(t)$ describes receipt if assignment were set to $t$.
\item $Y_i(t,d)$ allows assignment and receipt to affect outcomes separately.
\item The observed outcome is $Y_i=Y_i(T_i,D_i(T_i))$.
\item An exclusion restriction $Y_i(t,d)=Y_i(d)$ removes a direct effect of assignment on the outcome.
\end{items}
\EX{Section 3}
\end{frame}

\begin{frame}{Intention-to-treat effects}
Define the right-minus-left jump operator $\Delta_c$.
\[
\tau_Y=\Delta_c\E[Y_i\mid X_i=x],\qquad
\tau_D=\Delta_c\E[D_i\mid X_i=x].
\]
\begin{items}
\item $\tau_Y$ is the reduced form: the effect of assignment on the outcome under the relevant identification assumptions.
\item $\tau_D$ is the first stage: the change in treatment receipt induced by crossing the cutoff.
\item Assignment effects can be policy relevant even when the receipt effect needs further assumptions.
\end{items}
\EX{Section 3.2.1}
\end{frame}

\begin{frame}{Compliance types}
\begin{center}\renewcommand{\arraystretch}{1.5}
\begin{tabular}{@{}lcc@{}}\toprule
Type & $D_i(0)$ & $D_i(1)$\\\midrule
Never-taker & 0 & 0\\
Complier & 0 & 1\\
Always-taker & 1 & 1\\
Defier & 1 & 0\\\bottomrule
\end{tabular}\end{center}
Monotonicity, $D_i(1)\geq D_i(0)$, excludes defiers. Compliance type is latent because only one potential treatment is observed.
\EX{Section 3.2.2}
\end{frame}

\begin{frame}{The fuzzy RD local treatment effect}
Under continuity of the relevant potential quantities, exclusion, a nonzero first stage, and monotonicity,
\[
\boxed{\tau_{\mathrm{FRD}}=\frac{\tau_Y}{\tau_D}
=\E[Y_i(1)-Y_i(0)\mid D_i(1)>D_i(0),X_i=c].}
\]
\begin{items}
\item This is a local average treatment effect for cutoff compliers.
\item Always-takers and never-takers do not identify the receipt effect.
\item A treatment effect for all recipients is a different parameter.
\end{items}
\EX{Section 3.2.2}
\end{frame}

\begin{frame}{Estimation of a fuzzy RD effect}
\[
\widehat\tau_{\mathrm{FRD}}=\frac{\widehat\tau_Y}{\widehat\tau_D}.
\]
\begin{items}
\item The reduced form and first stage must use compatible samples and local specifications for this ratio calculation.
\item A joint fuzzy RD procedure selects bandwidths and accounts for covariance between numerator and denominator.
\item Dividing two separately optimized point estimates need not reproduce a joint fuzzy estimate.
\item Inference must account for smoothing bias in the ratio.
\end{items}
\EX{Sections 3.2.2--3.2.3}
\end{frame}

\begin{frame}{Fuzzy RD under local randomization}
\begin{items}
\item Invoke the local randomization conditions within a credible window.
\item Use assignment as an instrument for receipt within that window.
\item Exclusion and monotonicity yield a window complier average effect, under appropriate assignment assumptions.
\item Large sample TSLS inference and finite sample tests based on assignment require different arguments.
\end{items}
\EX{Sections 3.2.1--3.2.3}
\end{frame}

\begin{frame}{Weak assignment and validation}
\begin{items}
\item A small or imprecisely estimated first stage makes the ratio unstable.
\item Ordinary Wald intervals can fail under weak identification.
\item Report the first stage and use inference suited to weak assignment when necessary.
\item Covariate and density diagnostics concern the assignment rule. Exclusion and monotonicity need substantive arguments too.
\end{items}
\EX{Sections 3.2.4--3.2.5}
\end{frame}

\subsection{Discrete running variables}
\begin{frame}{A score with mass points}
\begin{items}
\item Many units can share a score when the running variable takes discrete values.
\item Count distinct score values overall and inside each estimation window.
\item Inspect the closest score values on both sides and the treatment rule at exact ties.
\item More observations at existing values improve precision there, but do not create observations closer to the cutoff.
\end{items}
\EX{Sections 4.1--4.2}
\end{frame}

\begin{frame}{The support gap around the cutoff}
Let $x_-$ and $x_+$ be the closest observed score values on each side.
\[
\E[Y_i\mid X_i=x_+]-\E[Y_i\mid X_i=x_-]
\]
is an observed cross-score contrast.
\begin{items}
\item If $x_-<c<x_+$ with a fixed gap, continuity alone does not identify the missing limits.
\item A cutoff effect then requires restrictions on how outcomes behave between support points.
\item A larger number of units at $x_-$ and $x_+$ does not eliminate this problem.
\end{items}
\EX{Section 4.3.1}
\end{frame}

\begin{frame}{Continuity methods with a discrete score}
\begin{items}
\item Local polynomial fits can be informative when there are many distinct values and the relevant smoothness assumptions are plausible.
\item Inference still depends on how extrapolation between mass points is justified.
\item Software adjustments for mass points improve practical bandwidth calculations.
\item Clustering by the score does not itself resolve the identification or approximation problem.
\end{items}
\EX{Section 4.3}
\end{frame}

\begin{frame}{Local randomization with a discrete score}
\begin{items}
\item Define a window containing a small, explicit set of score values.
\item State why assignment across those values can be treated as randomized.
\item Check predetermined covariates and assignment counts.
\item The estimand becomes an effect for the units in that window, with assumptions that differ from smooth extrapolation to a missing cutoff.
\end{items}
\EX{Section 4.4}
\end{frame}

\begin{frame}{Ties, selection, and interpretation}
\begin{items}
\item Whether equality triggers treatment must agree with the actual rule.
\item A tiny numerical shift can encode a tie rule, but it cannot supply new information.
\item Missing outcomes may themselves result from treatment, such as outcomes observed only for people who remain enrolled.
\item Identification of a score contrast does not automatically identify an effect for a stable population of observed survivors.
\end{items}
\EX{Sections 4.1 and 4.4}
\end{frame}

\subsection{Multidimensional RD designs}
\begin{frame}{Multiple cutoffs: two assignment structures}
\fig[.68]{cutoff_types.pdf}
\src{Noncumulative: each unit faces one group-specific cutoff. Cumulative: the same unit faces several thresholds, often for different treatment levels.}
\EX{Section 5.1.1, Figure 12}
\end{frame}

\begin{frame}{Cutoff-specific treatment effects}
For a noncumulative design, let $C_i$ denote the cutoff faced by unit $i$:
\[
T_i=\ind\{X_i\geq C_i\},\qquad
\tau(c)=\E[Y_i(1)-Y_i(0)\mid X_i=c,C_i=c].
\]
\begin{items}
\item Each cutoff has its own local comparison and identifying assumptions.
\item Differences between cutoff effects can reflect different populations and different score locations.
\item Comparisons across cutoffs require attention to uncertainty and dependence.
\end{items}
\EX{Section 5.1.2}
\end{frame}

\begin{frame}{Normalizing and pooling}
Define $\widetilde X_i=X_i-C_i$. All thresholds now occur at zero.
\[
\tau_{\mathrm{pool}}=\sum_{c\in\mathcal C}\omega(c)\tau(c).
\]
\begin{items}
\item Under suitable conditions, the normalized discontinuity averages cutoff effects.
\item The weights depend on cutoff-group frequencies and score densities at their thresholds.
\item They generally differ from equal weights and full-sample group shares.
\item Pooling obscures heterogeneity unless cutoff-specific effects are also shown.
\end{items}
\EX{Section 5.1.2}
\end{frame}

\begin{frame}{Extrapolation across cutoff groups}
\fig[.57]{cutoff_extrapolation.pdf}
\src{A constant untreated-outcome gap across score locations can link a lower-cutoff group's treated outcomes to a higher-cutoff group's controls.}
\EX{Section 5.1.2, Figure 14; \textcite{cktv2021}}
\end{frame}

\begin{frame}{The assumption behind extrapolation}
For cutoff groups $c_L<c_H$, define
\[
B(x)=\mu_{0,L}(x)-\mu_{0,H}(x).
\]
If $B(x)=B(c_L)$ for $x\in(c_L,c_H)$, then
\[
\tau_L(x)=\mu_{1,L}(x)-\mu_{0,H}(x)-B(c_L).
\]
\begin{items}
\item The lower-cutoff group is treated and the higher-cutoff group untreated at $x$.
\item The constant-gap condition supplies the missing counterfactual. It is an additional identification assumption.
\end{items}
\EX{Section 5.1.2}
\end{frame}

\begin{frame}{Multiple scores and a treatment boundary}
\fig[.51]{multiscore_regions.pdf}
\[
X_i=(X_{1i},X_{2i}),\qquad T_i=a(X_i).
\]
Treatment changes along a boundary $\mathcal B$, which contains many possible comparison points.
\EX{Section 5.2, Figure 16}
\end{frame}

\begin{frame}{Boundary point and pooled effects}
For a boundary point $b\in\mathcal B$, the target is
\[
\tau(b)=\E[Y_i(1)-Y_i(0)\mid X_i=b].
\]
\begin{items}
\item A signed distance to a particular point can implement a point-specific comparison under appropriate conditions.
\item Distance to the \emph{closest point anywhere on the boundary} pools different locations.
\item The pooled effect and point-specific effects answer different questions.
\item Distance units and the metric affect which observations are close.
\end{items}
\EX{Section 5.2.1}
\end{frame}

\begin{frame}{Geographic boundaries}
\begin{items}
\item Latitude and longitude jointly determine which side of a boundary a unit occupies.
\item Geographic proximity can coexist with sharp differences in sorting, services, or institutions.
\item Sparse settlement near the border increases extrapolation.
\item Spatial dependence and spillovers can affect inference and interpretation.
\item The causal claim requires a defensible treatment contrast at the chosen locations.
\end{items}
\EX{Section 5.2.3}
\end{frame}

\begin{frame}{Choosing the estimand and framework}
\small\renewcommand{\arraystretch}{1.45}
\begin{tabular}{@{}p{.26\linewidth}p{.64\linewidth}@{}}\toprule
Setting & Central question before estimation\\\midrule
Continuity, sharp RD & What effect do the outcome limits identify at $c$?\\
Local randomization & Which window and assignment mechanism are credible?\\
Fuzzy RD & What receipt effect is identified for which compliers?\\
Discrete score & What justifies the missing counterfactual across the support gap?\\
Multiple dimensions & Which cutoff group or boundary location defines the target?\\\bottomrule
\end{tabular}
\EX{Sections 2--6}
\end{frame}
